\documentclass[a4paper,12pt]{article}
\usepackage[italian]{babel}
\usepackage[utf8]{inputenc}
\usepackage[T1]{fontenc}
\usepackage{geometry}
\usepackage{setspace}
\geometry{margin=2.5cm}
\onehalfspacing

\title{Verifica scritta di Informatica - SQL e Algebra Relazionale}
\author{Prof. Fedeli Massimo}
\date{}

\begin{document}
	\maketitle
	
	\textbf{Cognome e Nome:} \underline{\hspace{8cm}} \hfill
	\textbf{Data:} \underline{\hspace{3cm}}
	
	\vspace{0.5cm}
	
	\textbf{Durata:} 2 ore \\
	\textbf{Obiettivo:} scrivere interrogazioni SQL e formulare interrogazioni in algebra relazionale su uno schema logico diverso rispetto a quelli già svolti in classe.
	
	\section*{Schema logico}
	
	\begin{flushleft}
		UTENTI(id\_utente PK, cognome, nome, classe, email, data\_iscrizione) \\
		LIBRI(id\_libro PK, titolo, autore, genere, anno\_pubblicazione, editore, copie\_totali) \\
		PRESTITI(id\_prestito PK, id\_utente FK, id\_libro FK, data\_prestito, data\_scadenza, data\_restituzione) \\
		PRENOTAZIONI(id\_prenotazione PK, id\_utente FK, id\_libro FK, data\_prenotazione, stato)
	\end{flushleft}
	
	\section*{Indicazioni}
	
	Scrivere le query richieste in linguaggio SQL. Quando richiesto, esprimere anche l'interrogazione in algebra relazionale. Evitare abbreviazioni non standard. Usare alias significativi.
	
	\section*{Quesiti}
	
	\begin{enumerate}
		\item Elencare cognome, nome e classe di tutti gli utenti iscritti alla biblioteca, ordinati per classe e cognome.
		
		\item Visualizzare titolo, autore e genere dei libri pubblicati dopo il 2018.
		
		\item Elencare i prestiti attualmente non restituiti, mostrando il nome dell'utente, il titolo del libro e la data di scadenza.
		
		\item Trovare tutti i libri del genere \texttt{Informatica} oppure \texttt{Matematica}, ordinati per autore.
		
		\item Contare quanti libri sono presenti per ciascun genere.
		
		\item Determinare, per ogni utente, il numero totale di prestiti effettuati. Mostrare anche gli utenti che non hanno mai effettuato prestiti.
		
		\item Elencare i titoli dei libri che non sono mai stati presi in prestito.
		
		\item Trovare gli utenti che hanno preso in prestito almeno 3 libri diversi.
		
		\item Visualizzare i libri che hanno un numero di prenotazioni maggiore di 2 con stato uguale a \texttt{attiva}.
		
		\item Trovare l'utente o gli utenti che hanno effettuato il maggior numero di prestiti.
		
		\item Elencare i libri attualmente disponibili, cioè quelli per cui il numero di copie in prestito al momento è minore del numero totale di copie possedute.
		
		\item Scrivere in algebra relazionale l'interrogazione che restituisce cognome e nome degli utenti che hanno almeno un prestito non ancora restituito.
		
		\item Scrivere in algebra relazionale l'interrogazione che restituisce i titoli dei libri prenotati almeno una volta.
		
		\item Trovare gli utenti che hanno effettuato almeno un prestito ma nessuna prenotazione.
		
		\item Per ciascun genere, calcolare l'anno medio di pubblicazione dei libri appartenenti a quel genere.
	\end{enumerate}
	
\end{document}
